\section{Background and Related Literature}
\label{sec:background}

\subsection{Institutional Background}

\paragraph{PDUFA (1992).}
The Prescription Drug User Fee Act, enacted in October 1992, authorized the FDA to collect fees from pharmaceutical manufacturers to supplement congressional appropriations for drug review \parencite{rep.dingellHR5952102ndCongress1992}. In exchange, the agency agreed to binding performance goals specifying median review times for new drug applications. The original PDUFA I (1993--1997) committed the FDA to reviewing standard applications within twelve months and priority applications within six months of receipt, with completion targets phased in over five years. Subsequent reauthorizations --- PDUFA II through VII --- have progressively tightened review-time targets, expanded the user-fee program to biologics, and added post-market safety obligations. The fee structure is renegotiated between the pharmaceutical industry and the FDA before each five-year reauthorization cycle, a feature that has long raised concerns about regulatory capture \parencite{darrowSpeedSafetyIndustry2017}.

PDUFA applies to applications for new molecular entities and new formulations submitted under the NDA pathway. It does not directly govern the Abbreviated New Drug Application (ANDA) pathway for generic drugs, which was created by the Hatch--Waxman Act of 1984 \parencite{rep.waxmanHR360598thCongress1984}. Until 2012, ANDA reviews were funded entirely by congressional appropriations, which contributed to a backlog of thousands of pending generic applications by the late 2000s.

\paragraph{GDUFA (2012).}
The Generic Drug User Fee Amendments, signed in July 2012, extended the user-fee model to generic drug reviews under the ANDA pathway \parencite{researchGDUFA2024}. GDUFA~I (2013--2017) collected fees from generic manufacturers and established performance goals for ANDA review times. A distinguishing feature of the initial GDUFA implementation was a substantial transition period (2013--2014) during which the FDA worked to clear a pre-existing backlog of approximately 2{,}500 unreviewed applications; the performance-goal era for new ANDA submissions effectively began in 2015. GDUFA~II (2018--2022) and GDUFA~III (2023--2027) have further refined the fee structure and review targets. Because 86.6\% of controlled-substance original approvals in this dataset came through the ANDA pathway, GDUFA is arguably the more directly relevant policy for this paper's central question.

\subsection{Theory of User-Fee Incentives}
\label{sec:theory}

The standard economic argument for why user-fee acceleration should change the composition of FDA approvals operates through the present value of approval. Let $V_j$ denote the present value of approval for drug $j$, $\pi_{jt}$ the per-period flow profit between approval and patent expiration, $\delta$ the discount factor, $T_j$ the patent expiration date, and $\tau_j$ the approval date. Then
\[
V_j \;=\; \sum_{t = \tau_j}^{T_j} \delta^{t - \tau_j}\,\pi_{jt}.
\]
A reduction in $\tau_j$ --- the approval date --- raises $V_j$ by extending the marketing window and front-loading the cash-flow stream. The marginal value of one additional period of marketing exclusivity, $\partial V_j / \partial \tau_j$, is increasing in $\pi_{jt}$. Holding patent term fixed, the firms with the most to gain from review acceleration are therefore those with the highest expected per-period profits during the on-patent window.

Three features of the controlled-substance segment plausibly place it above the cross-sectional mean of $\pi_{jt}$ during the relevant period. First, branded extended-release and abuse-deterrent opioid formulations command high list prices and have historically faced limited generic competition under patent protection \parencite{powellHowIncreasingMedical2020}. Second, demand for prescription opioid analgesics in the late twentieth and early twenty-first centuries grew rapidly relative to most non-controlled therapeutic categories, expanding the market in which these drugs competed. Third, the regulatory and reimbursement environment for chronic pain --- the largest indication for branded opioids --- supported broad insurance coverage and clinician adoption, raising effective per-patient revenue. Each of these features makes the theoretical prediction sharper rather than weaker: if the value of accelerated approval scales with expected on-patent profits, branded controlled substances should be near the top of the cross-sectional distribution of value gained from PDUFA.

The same logic, with the sign reversed, applies on the generic margin under GDUFA. Schedule~II generics carry meaningful additional manufacturing, distribution, and security costs --- DEA quota allocations, controlled storage, and Suspicious Order Monitoring requirements raise the fixed and variable costs of producing a controlled-substance generic relative to a non-controlled generic. If user-fee acceleration on the ANDA side reduces approval lags by an equal amount across applications, the resulting present-value boost is a smaller share of total project cost for non-controlled generics with low fixed costs than for controlled-substance generics, biasing the post-GDUFA composition shift toward non-CS ANDAs. Either way, the user-fee margin has a clear theoretical reason to move the controlled-substance share of approvals away from its pre-policy baseline.

A competing hypothesis holds that the operative mechanism behind PDUFA-era review-time improvements is not the demand-side shift created by user fees but the supply-side expansion in FDA review capacity that those fees financed. \textcite{ApprovalTimesNew} argue that the level of FDA staffing and review resources matters more than the source of funding. Under that view, PDUFA improved review-time performance primarily by enabling the FDA to hire and retain expertise, and any composition effect would operate through the supply of review capacity rather than through demand-side incentives for manufacturers. Because both channels were activated simultaneously by PDUFA, the reduced-form approval-side test in this paper does not separate them; what it does is bound the joint effect that any plausible mechanism could have produced through the approval-composition margin.

\subsection{Related Empirical Literature}

\paragraph{PDUFA, review speed, and approval.}
The most direct predecessor is \textcite{berndtIndustryFundingFDA2005}, who use drug-level panel data to show that PDUFA substantially reduced NDA review times and increased approval rates while finding little evidence of increased post-approval withdrawals. \textcite{olsonPDUFAInitialUS2009} examines launch decisions and finds that PDUFA increased the probability of U.S.\ initial launch conditional on development, consistent with the thesis that faster review improves the expected return to development. \textcite{olsonFirmCharacteristicsSpeed1997} documents substantial heterogeneity in pre-PDUFA review times across firm types, suggesting that any compositional effects of user-fee acceleration may depend on which firms were most constrained by review delays before 1992. \textcite{dranoveImportantDrugsReach1994} provide an early analysis of whether important drugs reach the market faster, finding limited evidence of differential queue position by therapeutic value.

\paragraph{Post-market safety.}
\textcite{carpenterDrugreviewDeadlinesSafety2008} exploit the fact that drugs approved just before PDUFA deadlines received less review time than those approved just after, finding that proximity-to-deadline approvals are associated with higher rates of subsequent safety-related labeling changes. \textcite{darrowFDAApprovalRegulation2020} review the broader evidence and argue that expedited pathways have reduced clinical evidence requirements at the time of approval, raising concerns about downstream safety for certain product classes. Both papers are informative about the speed--safety margin but do not address the market-composition question examined here.

\paragraph{Welfare analysis.}
\textcite{philipsonCostbenefitAnalysisFDA2008} develop a formal welfare model of the FDA approval decision and use it to estimate that approval delays under the pre-PDUFA regime imposed large welfare costs. The framework emphasizes the consumer-surplus value of earlier access to beneficial drugs but does not model externalities associated with abuse potential or diversion. That omission is consequential for controlled substances, where the social cost of approval may include downstream misuse spillovers that are not internalized by approving firms or patients.

\paragraph{Generic drugs and GDUFA.}
\textcite{berndtGenericDrugUser2018} provide the first systematic economic analysis of GDUFA, documenting that the reform reduced ANDA review times and approval lags. They also analyze barriers to generic entry and argue that regulatory delays were a significant constraint on competition in the pre-GDUFA era. Their analysis is primarily about generic market structure and does not focus on controlled-substance composition.

\paragraph{Controlled substances and opioid supply.}
\textcite{alpertSupplySideDrugPolicy2018a} show that the introduction of abuse-deterrent OxyContin in 2010 reduced prescription opioid misuse but was followed by substitution to heroin, illustrating supply-side spillovers in controlled-substance markets. \textcite{powellEVOLVINGCONSEQUENCESOXYCONTIN2021} extend this analysis through the fentanyl era. Neither paper studies the FDA approval side of the supply chain. This paper contributes by examining whether the upstream regulatory decision environment shaped the composition of what entered the market.

\paragraph{Gap.}
No existing study examines whether PDUFA or GDUFA changed the controlled-substance share of FDA approvals. The literature has studied review times, launch timing, post-market safety, and welfare under PDUFA, and market structure under GDUFA; the composition question has not been addressed. This paper fills that gap and reports a precisely estimated null whose magnitude bounds, rather than merely fails to reject, the user-fee composition hypothesis.
